\begin{abstract}
Let $X$ be a finite connected graph and let a finite group $R$ act regularly
on a receipt fiber over every vertex. We classify all regular $R$-covariant
graph lifts obtained from inverse-compatible edge receipts. A voltage
assignment $\alpha$ defines a path-receipt homomorphism
\[
\rho_{\alpha}:\pi_1(X,r)\longrightarrow R.
\]
Vertex gauge conjugates this representation, and a spanning-tree gauge
normalizes every tree-edge receipt to the identity. Consequently, gauge
classes of regular lifts are in bijection with $R$-conjugacy classes of
representations of the graph fundamental group.

If $H_{\alpha}=\im(\rho_{\alpha})$, then the lifted graph has exactly
$[R:H_{\alpha}]$ connected components, each containing
$|V(X)|\,|H_{\alpha}|$ vertices. In particular, the lift is connected if and
only if $H_{\alpha}=R$. The finite cycle-multiplication theorem is recovered
when $\pi_1(X)$ has rank one: then $H_{\alpha}=\langle h\rangle$, and receipt
order determines both component count and complete return length.

The paper packages the theorem with an executable finite census and an exact
raw-graph binding protocol. The protocol separates the abstract algebra from
the additional work required to derive a receipt group, regular fiber action,
and edge receipts from an unlabeled graph. The mathematical results belong to
the standard voltage-graph and regular-cover setting; no physical realization
or novelty claim over that theory is made.
\end{abstract}
